% Exemplo de Uso de Figuras na Dissertação UNICAMP-FT
% Este arquivo demonstra as melhores práticas para inclusão de figuras

% ============================================================================
% EXEMPLOS POR TIPO DE FIGURA
% ============================================================================

% ----------------------------------------------------------------------------
% 1. FIGURA SIMPLES (Capítulo)
% ----------------------------------------------------------------------------
\begin{figure}[htbp]
    \centering
    \includegraphics[width=0.8\textwidth]{figuras/capitulos/cap01-introducao/fig-cap01-01-arquitetura-hiperespectral}
    \caption{Arquitetura geral de um sistema de imageamento hiperespectral embarcado, 
             mostrando os componentes principais: sensor, processamento e armazenamento.}
    \label{fig:cap01-arquitetura}
\end{figure}

% Referência no texto:
% Como mostrado na \autoref{fig:cap01-arquitetura}, a arquitetura proposta...

% ----------------------------------------------------------------------------
% 2. DIAGRAMA DE BLOCOS
% ----------------------------------------------------------------------------
\begin{figure}[htbp]
    \centering
    \includegraphics[width=\textwidth]{figuras/diagramas/diag-arquitetura-fpga-gpu}
    \caption{Diagrama de blocos da arquitetura heterogênea FPGA-GPU proposta 
             para processamento hiperespectral em tempo real.}
    \label{fig:diag-fpga-gpu}
\end{figure}

% ----------------------------------------------------------------------------
% 3. GRÁFICO DE PERFORMANCE
% ----------------------------------------------------------------------------
\begin{figure}[htbp]
    \centering
    \includegraphics[width=0.9\textwidth]{figuras/graficos/graf-performance-comparativa}
    \caption{Comparação de performance entre diferentes arquiteturas para 
             processamento hiperespectral: throughput (fps) versus consumo 
             energético (W) para datasets de referência.}
    \label{fig:graf-performance}
\end{figure}

% ----------------------------------------------------------------------------
% 4. ESQUEMA TÉCNICO
% ----------------------------------------------------------------------------
\begin{figure}[htbp]
    \centering
    \includegraphics[width=0.7\textwidth]{figuras/esquemas/esq-sistema-embarcado}
    \caption{Esquema detalhado do sistema embarcado proposto, incluindo 
             hierarquia de memória, unidades de processamento e interfaces 
             de comunicação.}
    \label{fig:esq-sistema}
\end{figure}

% ----------------------------------------------------------------------------
% 5. DUAS FIGURAS LADO A LADO
% ----------------------------------------------------------------------------
\begin{figure}[htbp]
    \centering
    \begin{subfigure}[b]{0.48\textwidth}
        \centering
        \includegraphics[width=\textwidth]{figuras/graficos/graf-throughput}
        \caption{Throughput}
        \label{fig:throughput}
    \end{subfigure}
    \hfill
    \begin{subfigure}[b]{0.48\textwidth}
        \centering
        \includegraphics[width=\textwidth]{figuras/graficos/graf-latencia}
        \caption{Latência}
        \label{fig:latencia}
    \end{subfigure}
    \caption{Análise de performance: (a) throughput médio e (b) latência 
             para diferentes configurações de hardware.}
    \label{fig:performance-analise}
\end{figure}

% Referência: A \autoref{fig:throughput} mostra que...

% ----------------------------------------------------------------------------
% 6. FIGURA COM MÚLTIPLAS PARTES (2x2)
% ----------------------------------------------------------------------------
\begin{figure}[htbp]
    \centering
    \begin{subfigure}[b]{0.48\textwidth}
        \centering
        \includegraphics[width=\textwidth]{figuras/capitulos/cap04-resultados/fig-cap04-01-fpga-results}
        \caption{Resultados FPGA}
        \label{fig:results-fpga}
    \end{subfigure}
    \hfill
    \begin{subfigure}[b]{0.48\textwidth}
        \centering
        \includegraphics[width=\textwidth]{figuras/capitulos/cap04-resultados/fig-cap04-02-gpu-results}
        \caption{Resultados GPU}
        \label{fig:results-gpu}
    \end{subfigure}
    
    \vspace{0.5cm}
    
    \begin{subfigure}[b]{0.48\textwidth}
        \centering
        \includegraphics[width=\textwidth]{figuras/capitulos/cap04-resultados/fig-cap04-03-hybrid-results}
        \caption{Resultados Híbridos}
        \label{fig:results-hybrid}
    \end{subfigure}
    \hfill
    \begin{subfigure}[b]{0.48\textwidth}
        \centering
        \includegraphics[width=\textwidth]{figuras/capitulos/cap04-resultados/fig-cap04-04-comparison}
        \caption{Comparação}
        \label{fig:results-comparison}
    \end{subfigure}
    
    \caption{Resultados experimentais para diferentes arquiteturas: 
             (a) FPGA standalone, (b) GPU standalone, (c) arquitetura 
             híbrida proposta, e (d) comparação geral de performance.}
    \label{fig:resultados-completos}
\end{figure}

% ----------------------------------------------------------------------------
% 7. TABELA COMO FIGURA (para tabelas complexas)
% ----------------------------------------------------------------------------
\begin{figure}[htbp]
    \centering
    \includegraphics[width=\textwidth]{figuras/tabelas/tab-comparacao-tecnicas}
    \caption{Comparação detalhada entre técnicas de processamento hiperespectral 
             embarcado, incluindo métricas de performance, consumo energético 
             e precisão de classificação.}
    \label{fig:tab-comparacao}
\end{figure}

% ============================================================================
% COMANDOS ÚTEIS PARA FIGURAS
% ============================================================================

% ----------------------------------------------------------------------------
% Ajustes de Tamanho
% ----------------------------------------------------------------------------
% \includegraphics[width=0.5\textwidth]{figura}     % 50% da largura do texto
% \includegraphics[width=\textwidth]{figura}        % 100% da largura do texto
% \includegraphics[height=5cm]{figura}              % Altura específica
% \includegraphics[scale=0.8]{figura}               % Escala relativa
% \includegraphics[width=10cm,height=8cm]{figura}   % Dimensões específicas

% ----------------------------------------------------------------------------
% Posicionamento
% ----------------------------------------------------------------------------
% [h]    - here (aqui)
% [t]    - top (topo da página)
% [b]    - bottom (final da página)
% [p]    - page (página separada)
% [htbp] - hierarquia de preferências (recomendado)
% [H]    - força posição exata (requer \usepackage{float})

% ----------------------------------------------------------------------------
% Rotação
% ----------------------------------------------------------------------------
% \includegraphics[angle=90,width=\textwidth]{figura}  % Rotacionar 90°

% ----------------------------------------------------------------------------
% Figuras Lado a Lado (alternativa simples)
% ----------------------------------------------------------------------------
% \begin{figure}[htbp]
%     \centering
%     \includegraphics[width=0.45\textwidth]{figura1}
%     \hfill
%     \includegraphics[width=0.45\textwidth]{figura2}
%     \caption{Duas figuras lado a lado}
%     \label{fig:lado-a-lado}
% \end{figure}

% ============================================================================
% PACOTES NECESSÁRIOS NO PREÂMBULO
% ============================================================================
% \usepackage{graphicx}      % Para \includegraphics
% \usepackage{subcaption}    % Para subfiguras
% \usepackage{float}         % Para posicionamento [H]
% \usepackage{caption}       % Para personalizar captions

% ============================================================================
% CONFIGURAÇÕES RECOMENDADAS NO PREÂMBULO
% ============================================================================
% % Caminho padrão para figuras
% \graphicspath{{figuras/}}
% 
% % Extensões de arquivo (ordem de prioridade)
% \DeclareGraphicsExtensions{.pdf,.png,.jpg,.eps}
% 
% % Configuração de captions
% \captionsetup{
%     font=small,
%     labelfont=bf,
%     textfont=it,
%     justification=centering
% }

% ============================================================================
% DICAS IMPORTANTES
% ============================================================================
% 1. Sempre usar \autoref{} para referenciar figuras
% 2. Colocar figuras próximas ao texto que as referencia
% 3. Captions devem ser autoexplicativas
% 4. Usar labels descritivos e únicos
% 5. Preferir formato PDF para figuras vetoriais
% 6. Manter resolução mínima de 300 DPI para figuras rasterizadas
% 7. Testar visualização em diferentes tamanhos
% 8. Verificar legibilidade quando impressa em preto e branco
